% !PRINTED 600
\weberpage{600}
In diesem Körper sei $A$ eine Kongruenzgruppe mit dem Modul
$M$ und dem Exkludenten $S$. Die Gruppe $A$ enthält, wie wir gesehen
 haben, unendlich viele \gesperrt{ganze} Zahlen. Wir fragen nach
der Anzahl $T(t)$ der ganzen Zahlen in $A$, die durch irgend ein
gegebenes ganzes zu $M$ teilerfremdes Ideal $\mathfrak{m}$ teilbar sind, deren
Norm nicht größer als die positive Zahl $t$ ist.

Es sei $(\alpha_1, \alpha_2)$ eine Basis von $\mathfrak{m}$, also nach \S\,91, (11):
\weq{1}{\alpha_1 = a_{22}\,a, \quad \alpha_2 = a_{22}\,\dfrac{b + \sqrt{\Delta}}{2}.}

Hierin bedeute $a_{22}$ die größte in $\mathfrak{m}$ aufgehende ganze rationale
Zahl, $a_{22}\,a$ die kleinste durch $\mathfrak{m}$ teilbare ganze rationale Zahl;
$b$ ist eine Wurzel der Kongruenz:
\[
b^2 \equiv \Delta \pmod{4a},
\]
und es ist:
\weq{2}{\Delta = b^2 - 4ac.}
\weq{3}{N(\mathfrak{m}) = a_{22}^2\,a.}
Da $\mathfrak{m}$ relativ prim zu $M$ vorausgesetzt ist, so können wir nach
\S\,161, (6) in jeder in $A$ vorkommenden Zahlklasse nach $M$ eine
durch $\mathfrak{m}$ teilbare Zahl $\alpha_0$ bestimmen. Diese Zahl hat die Form
\[
\alpha_0 = x_0\alpha_1 + y_0\alpha_2,
\]
und wenn wir
\weq{4}{\begin{aligned} \alpha &= \alpha_0 + M(x\alpha_1 + y\alpha_2) \\ &= (x_0 + Mx)\alpha_1 + (y_0 + My)\alpha_2 \end{aligned}}
setzen und $x, y$ alle positiven und negativen ganzen rationalen
Zahlen durchlaufen lassen, so durchläuft $\alpha$ die Gesamtheit der
durch $\mathfrak{m}$ teilbaren und nach $M$ mit $\alpha_0$ kongruenten ganzen Zahlen
der Gruppe $A$.

Um alle durch $\mathfrak{m}$ teilbaren ganzen Zahlen in $A$ zu erhalten,
haben wir ein vollständiges Restsystem nach $M$ in $\Omega$ zu suchen
und darunter die in $A$ enthaltenen durch $\mathfrak{m}$ teilbaren Zahlen $\alpha_0$
auszuwählen.

Soll die Norm dieser Zahl $\alpha$ nicht größer als $t$ sein, so ergibt
 sich dafür, wenn wir die quadratische Form
\[
a x^2 + b x y + c y^2
\]
mit $\varphi(x, y)$ bezeichnen, die Bedingung:
\[
\varphi(x_0 + Mx, \; y_0 + My) \leqq \dfrac{t}{a_{22}^2\,a}, \qquad [\S\,91, (2), (4)].
\]
